\bookchapter{第四十七章：速查手册与资源}

\chapterintro{本章是一份可以独立使用的速查手册。当你记不清某个原则或需要快速决策时，直接翻到这里。建议你把它当成“口袋卡片”：五大原则速查、工具决策树、常见问题、学习资源，一页不落。}

\booksection{一、思维模型速查卡}

\begin{center}
\begin{tabularx}{\linewidth}{|X|X|X|}
\hline
\textbf{原则} & \textbf{一句话} & \textbf{操作要点} \tabularnewline
\hline
Think Before Coding & 先想清楚再开口 & 回答三问：做什么？结果长啥样？有什么约束？ \tabularnewline
Simplicity First & 用最少代码解决问题 & 每次只做一个功能，砍掉不必要的复杂度 \tabularnewline
Surgical Changes & 精准修改不越界 & 明确指定修改范围，拒绝AI的"顺手改一下" \tabularnewline
Goal-Driven Execution & 说"做什么"不说"怎么做" & 给目标不给步骤，让AI自主选择路径 \tabularnewline
Iterate \& Verify & 生成-验证-修正循环 & 三层验证：功能、边界、架构 \tabularnewline
\hline
\end{tabularx}
\end{center}

\begin{center}\small\color{gray}五大思维原则速查卡\end{center}

\booksection{二、工具选择决策树}

当你需要选择 AI 编程工具时，按照以下决策路径选择。注意：这不是“哪个最好”的排名，而是“哪个最适合你”的判断框架。

你要做什么？

零基础，想快速体验 选云端 Builder (v0/Bolt/Replit)

日常开发，需要代码补全 选 IDE 内置 AI (Cursor/Trae)

大型项目，需要批量操作 选终端 Agent (Claude Code)

快速原型验证想法 选云端 Builder 或 Trae SOLO

\begin{center}\small\color{gray}图 6：AI 编程工具选择决策树\end{center}

\booksection{三、常见问题 FAQ}

\textbf{Q1：我完全不会编程，能用 AI 编程吗？}

能。2026 年的 AI 编程工具已经足够强大，你完全可以用自然语言描述需求，让 AI 生成代码。但有两点建议：第一，从云端 Builder 工具（如 v0、Bolt）开始，它们不需要配置开发环境；第二，边用边学基础概念——不需要精通语法，但你需要能“看懂”AI 写的代码的大意，这样才能有效验收。

\textbf{Q2：我应该学哪个编程语言？}

取决于你要做什么。做网站前端——JavaScript/TypeScript；做数据处理或 AI——Python；做企业级后端——Java 或 Go。但更重要的是：在 AI 编程时代，“语言”的重要性在下降。AI 能帮你处理语法细节，你需要理解的是“程序的结构”——什么是函数、什么是循环、什么是数据结构。这些概念是跨语言的，一旦理解，换什么语言都快。

\textbf{Q3：AI 写的代码有版权问题吗？}

截至 2026 年 8 月，这仍是一个有争议的法律问题。一般来说，AI 生成的代码在大多数司法管辖区不被视为有明确的版权所有者。但如果你用于商业项目，建议：(1) 了解你使用的 AI 工具的服务条款中关于知识产权的条款；(2) 对关键代码进行人工审查和修改，使其具有你的创造性投入；(3) 避免让 AI 直接复制可能有版权的开源代码。

\textbf{Q4：AI 会取代程序员吗？}

短期内不会。AI 取代的不是“程序员”这个职业，而是“逐行编码”这项工作。就像计算器没有取代数学家，而是让数学家从繁琐的手算中解放出来一样。AI 编程工具让程序员从“写代码”转向“想问题”——这实际上提高了程序员的价值上限。未来的竞争力不在于“能写多少行代码”，而在于“能解决多复杂的问题”。

\textbf{Q5：如何判断 AI 的产出是否正确？}

三层验证法。第一层：功能验证——它能不能跑？输入输出对不对？这是最基本的。第二层：边界验证——空输入、超长输入、异常输入会怎样？这能发现隐藏的问题。第三层：架构验证——代码结构是否合理？有没有硬编码？是否便于扩展？这需要你有一定的技术判断力。如果你无法判断第三层，至少做到前两层——能跑、边界正常，就比大多数初学者强了。

\booksection{四、推荐学习资源}

以下是截至 2026 年 8 月的推荐学习资源。请注意，互联网资源链接可能变化，如果某个链接失效，搜索其标题通常能找到最新版本。

\begin{center}
\begin{tabularx}{\linewidth}{|X|X|X|}
\hline
\textbf{类型} & \textbf{资源} & \textbf{推荐理由} \tabularnewline
\hline
原始文献 & Andrej Karpathy 的 CLAUDE.md & 五大原则的源头，70行但字字珠玑 \tabularnewline
原始文献 & Anthropic 官方文档 (Claude Code) & 终端Agent的最佳实践指南 \tabularnewline
教程 & Vibe Coding 社区整理的 Skills 合集 & 社区实践者总结的实战技巧集合 \tabularnewline
教程 & Cursor 官方文档和 YouTube 频道 & IDE内置AI的权威使用指南 \tabularnewline
社区 & GitHub 上的 awesome-claude-code & Claude Code 的最佳实践和插件集 \tabularnewline
社区 & Reddit 的 r/cursor 社区 & 全球Cursor用户的经验分享和讨论 \tabularnewline
基础 & MDN Web Docs (developer.mozilla.org) & 前端开发最权威的参考文档 \tabularnewline
基础 & Python 官方教程 (docs.python.org) & Python学习的官方起点 \tabularnewline
\hline
\end{tabularx}
\end{center}

\begin{center}\small\color{gray}推荐学习资源清单\end{center}

\booksection{五、结语：工具在手，方法论在心}

回到本书的核心信念：\textbf{工具每三个月迭代一次，方法论每三年才进化一次}。在 2026 年 8 月的今天，AI 编程工具的能力已经强大到令人惊叹——你能用一句话创建一个网站，用一段描述生成一个全栈应用，用一条指令完成一个复杂重构。但工具越强大，方法论的重要性反而越突出——因为工具的强大放大了你的每一个决策。

如果你的需求清晰，AI 能帮你高效实现；如果你的需求模糊，AI 也会高效地做出错误的东西。如果你遵循简单至上，AI 能帮你快速验证想法；如果你贪多求全，AI 也会帮你快速堆积技术债务。如果你精准修改，AI 能帮你外科手术般地优化代码；如果你指令模糊，AI 也会帮你大刀阔斧地引入 bug。

这就是为什么本书花了大部分篇幅讲方法论而非工具操作。不是因为工具不重要，而是因为\textbf{方法论是工具的“乘数”}——方法论越好，工具的价值越大；方法论为零，工具再强也只是零乘以任何数。

你不再是一个人在编程。你是一个指挥官，AI 是你的军团。你的思考质量决定了军团的战斗力。从今天开始，用方法论武装自己的思维，用工具放大自己的能力——这就是 AI 编程的全部秘密。

\booksection{参考文献}

\begin{enumerate}
  \item Andrej Karpathy. (2025). Vibe Coding. Social media post introducing the concept of vibe-driven programming with AI.
  \item Andrej Karpathy. (2025). CLAUDE.md - Project guidelines for AI-assisted development. GitHub repository, 70+K stars.
  \item Anthropic. (2026). Claude Code: Terminal-native AI programming. Official documentation and best practices.
  \item Cursor Team. (2026). Cursor 3: The next generation of AI-first IDE. Release notes and documentation.
  \item OpenAI. (2026). Codex CLI: Bringing AI to the terminal. Technical documentation.
  \item Model Context Protocol (MCP) Specification. (2025-2026). Open protocol for AI-tool integration.
  \item Vibe Coding Community. (2026). Community-curated skills and best practices for AI-assisted programming.
  \item Spec-Driven Development (SDD) methodology. (2026). Emerging practices in specification-first AI programming.
\end{enumerate}
